\begin{afterword}
\summaryhead

The post accepts the standard analysis of the falling-tree riddle: one side means acoustic vibrations by ``sound,'' the other means auditory experience, so the dispute is about a word. It asks why people argue about it at all.

The post returns to the author's invented ``bleggs'' and ``rubes,'' objects sorted by color, shape and contents. Once every feature of an object is known, the question ``is it really a blegg?'' has nothing left to stand for. The author compares two neural networks. In the first, every unit is an observable feature. In the second, the features connect through a central category unit, which can stay active or inactive after all features are fixed. The suggestion is that the brain resembles the second network, and that this dangling unit is what the leftover question feels like from the inside.

The same feeling, the post argues, drives disputes about the tree and about whether Pluto is a planet. People see their intuitions not as intuitions but as the world itself, and so they reject explanations of how their thinking goes wrong.

\responsehead

The post has a clear question and a clear answer. The question is why people go on arguing about the falling tree after the two meanings of ``sound'' have been pointed out. The analysis into two meanings is old; the post calls it ``the standard rationalist view'' and claims nothing for it. The answer is the post's own: brains categorize through a central unit, and the unit's leftover activation is what the leftover question feels like.

That answer is a hypothesis about brains, and the post gives it no evidence about brains. Network 2 is offered as ``a far better candidate for being something vaguely like how the human brain works,'' three hedges in one clause, and the reasons given are that it is ``fast, cheap, scalable.'' Those are reasons an engineer might choose it. They are not observations that brains use it. No study of categorization, human or animal, is cited. The diagram makes the explanation easy to see, but seeing it is not evidence for it. A reader who finishes the post should hold its account of the tree, of Pluto and of stubborn disagreement as conditional on a guess that the post does not test.

The post also states a disputed theory of perception as fact. When you look at a green cup, it says, you are seeing ``a picture reconstructed in your visual cortex,'' and ``that is what you are seeing.'' Philosophers of perception have argued about exactly this for centuries. Direct realists hold that we see the cup, and that the picture in the cortex is how we see it, not what we see. The post settles the question in a dash and then uses it as the model for how intuitions work.

The ending turns the account toward the author's critics. People who reject ``everything you try to say about how the native cognitive algorithm goes astray'' are said to do so because they cannot see their intuitions as intuitions. The link on ``cling to their intuitions'' goes to the author's reply to dissenting commenters. My reading, from that link, is that the critics the post had in mind included the author's own. The post offers no test that would separate a reader who clings to an intuition from one who has a sound objection.

In short: a clear question answered with an untested guess about brain design, a contested theory of perception stated as fact, and an ending that files disagreement under the bias it describes.
\end{afterword}
